\section{Worked example: randomness versus interference}

Classical random choice and quantum superposition can produce identical
single-measurement frequencies. They are not the same theory. This example
uses an intermediate phase change to distinguish them.

\subsection{Classical fair coin model}

Imagine a classical variable R chosen by a fair coin:
\[
P(R=0)=\frac12,\qquad P(R=1)=\frac12.
\]
Once chosen, R is one definite value even if the observer has not looked.
Applying classical NOT changes 0 to 1 and 1 to 0, but the distribution remains
half and half. There is no operation that makes the probability of one hidden
branch cancel the other branch.

\subsection{Quantum preparation with Hadamard}

\begin{lstlisting}
MAIN-PROCESS QuantumCoin
CREATETOKEN -I Q
SET Q 0p
H -I Q -O Q
MEASURE -I Q
RETURNVALS Q
\end{lstlisting}
Before measurement,
\[
\ket0\xrightarrow{H}
\ket+=\frac{\ket0+\ket1}{\sqrt2}.
\]
The amplitudes are $1/\sqrt2$ and $1/\sqrt2$; squaring magnitudes gives the
same half-and-half distribution as the classical coin. If measurement happens
now, this experiment alone cannot distinguish the two descriptions.

\subsection{Insert phase and recombine paths}

\begin{lstlisting}
MAIN-PROCESS PhaseInterference
CREATETOKEN -I Q
SET Q 0p
H -I Q -O Q
Z -I Q -O Q
H -I Q -O Q
MEASURE -I Q
RETURNVALS Q
\end{lstlisting}
Track the state:
\[
\ket0
\xrightarrow{H}
\frac{\ket0+\ket1}{\sqrt2}
\xrightarrow{Z}
\frac{\ket0-\ket1}{\sqrt2}
\xrightarrow{H}
\ket1.
\]

\subsection{Term-by-term explanation}

\begin{description}
  \item[Superposition] After the first H, both basis states have coherent
  amplitudes. No measurement has selected one.
  \item[Phase] Z multiplies only the $\ket1$ amplitude by $-1=e^{i\pi}$. Its
  magnitude remains $1/\sqrt2$, so immediate 0/1 probabilities remain equal.
  \item[Relative phase] The plus relationship between amplitudes becomes a
  minus relationship. A global minus on both terms would not matter; changing
  only one term does.
  \item[Interference] The final H mixes the basis amplitudes. Contributions to
  the $\ket0$ output cancel, while contributions to $\ket1$ add.
  \item[Measurement] The final state is $\ket1$, so measurement returns 1 with
  probability one.
\end{description}

The cancellation can be seen directly from matrix multiplication:
\[
H\frac1{\sqrt2}
\begin{bmatrix}1\\-1\end{bmatrix}
=
\frac12
\begin{bmatrix}1&1\\1&-1\end{bmatrix}
\begin{bmatrix}1\\-1\end{bmatrix}
=
\frac12\begin{bmatrix}0\\2\end{bmatrix}
=\begin{bmatrix}0\\1\end{bmatrix}.
\]
Classical probability weights are nonnegative and would add rather than cancel.
Quantum amplitudes are signed or complex, so phase can create destructive and
constructive interference before the Born rule turns amplitudes into
probabilities.

\subsection{Comparison table}

\begin{center}
\begin{tabularx}{\textwidth}{>{\bfseries}lXX}
\toprule
Question & Classical fair coin & H--Z--H quantum circuit\\
\midrule
State before observation & One hidden definite bit sampled fairly &
Coherent complex amplitude vector\\
Effect of middle sign/phase & No probability analogue that preserves a fair
distribution and later cancels a branch & Z changes relative phase without
changing immediate basis probabilities\\
Can paths cancel? & No; probabilities add & Yes; amplitudes add and may cancel\\
Final result in this example & Still a distribution unless another classical
rule forces a value & Deterministic 1\\
\bottomrule
\end{tabularx}
\end{center}

